% (c) Nikita Lisitsa, lisyarus@gmail.com, 2026

\documentclass[10pt]{beamer}

\usepackage[T2A]{fontenc}
\usepackage[russian]{babel}
\usepackage{minted}

\usepackage[most]{tcolorbox}
\tcbuselibrary{minted}

\usepackage{graphicx}
\graphicspath{ {./images/} }

\usetheme{metropolis}

\usepackage{adjustbox}

\usepackage{color}
\usepackage{soul}

\usepackage{hyperref}

\definecolor{codebg}{RGB}{29,35,49}

\setminted{fontsize=\footnotesize}

\setlength\partopsep{-\topsep}

\definecolor{blue}{rgb}{0,0,1}
\definecolor{red}{rgb}{1,0,0}

\usemintedstyle{tango}

\makeatletter
\newcommand{\slideimage}[1]{
  \begin{figure}
    \begin{adjustbox}{width=\textwidth, totalheight=\textheight-2\baselineskip-2\baselineskip,keepaspectratio}
      \includegraphics{#1}
    \end{adjustbox}
  \end{figure}
}
\makeatother

\title{Компьютерная графика}
\subtitle{Практика 9: Variance shadow mapping}
\date{2026}

\setbeamertemplate{footline}[frame number]

\begin{document}

\frame{\titlepage}

\begin{frame}[fragile]
\frametitle{Практика 9}
\slideimage{0.png}
\end{frame}

\begin{frame}[fragile]
\frametitle{Задание 1}
Подбираем проекцию для shadow map, в которую попадает вся сцена
\begin{itemize}
\item При загрузке 3D модели вычисляем её bounding box (на C++ будет полезен файл \mintinline{cpp}|<math/box.hpp>|)
\item За центр проекции берём центр bounding box'а
\item При рендеринге, по 8-ми вершинам bounding box'а находим размер области проекции по \(X\) (в коде это \mintinline{cpp}|light_X|) как \(\max\limits_V |(V-C)\cdot X|\)
\item Аналогично по \(Y\) и \(Z\)
\item Строим матрицу ортографической проекции, в shadow map теперь должна попасть вся сцена
\end{itemize}
\end{frame}

\begin{frame}[fragile]
\frametitle{Задание 1}
\slideimage{1.png}
\end{frame}

\begin{frame}[fragile]
\frametitle{Задание 2}
Создаём текстуру для VSM
\begin{itemize}
\item Создаём ещё одну текстуру shadow map с форматом \mintinline{cpp}|RG32F|, создаём для неё texture view, добавляем её в lights bind group
\item При генерации shadow map (первый render pass) добавляем эту текстуру в качестве color target
\item В shadow pipeline добавляем фрагментный шейдер который выводит цвет \mintinline{wgsl}|vec4f(z, z * z, 0.0, 0.0)|
\item В дебажном выводе shadow map теперь рисуем значение из этой текстуры (нужно добавить соответствующий binding в шейдер)
\item \textbf{NB:} Она будет желтоватой, так как в ней нулевой синий канал
\end{itemize}
\end{frame}

\begin{frame}[fragile]
\frametitle{Задание 2}
\slideimage{2.png}
\end{frame}

\begin{frame}[fragile]
\frametitle{Задание 3}
Реализуем VSM через неравенство Чебышева
\begin{itemize}
\item Меняем наш shadow sampler на линейную фильтрацию (min/mag фильтры)
\item В основной шейдер добавляем binding для нашей новой текстуры
\item Используем её вместо старой текстуры глубины для реализации variance shadow map (см. слайды лекции): в текстуре записаны \(\mu = \mathbb E [z]\) и \(\mathbb E [z^2]\), по ним нужно вычислить \(\sigma^2\) и финальное значение освещённости
\end{itemize}
\end{frame}

\begin{frame}[fragile]
\frametitle{Задание 3}
\slideimage{3.png}
\end{frame}

\begin{frame}[fragile]
\frametitle{Задание 4}
Добавляем shadow bias через shader derivatives
\begin{itemize}
\item При генерации shadow map, используя \mintinline{wgsl}|dpdx()| и \mintinline{wgsl}|dpdy()| вычисляем поправку к \(\mathbb E[z^2]\) (см. слайды лекции)
\item \textbf{NB:} Скорее всего, картинка почти не изменится
\end{itemize}
\end{frame}

\begin{frame}[fragile]
\frametitle{Задание 4}
\slideimage{3.png}
\end{frame}

\begin{frame}[fragile]
\frametitle{Задание 5}
Исправляем light bleeding
\begin{itemize}
\item Подбираем значение \(\delta\) и переводим диапазон значений освещённости (посчитанных через неравенство Чебышева):
\begin{itemize}
\item Значение \([0, \delta]\) должны перейти в \(0\)
\item Значение \([\delta, 1]\) должны перейти в \([0, 1]\)
\end{itemize}
\item Слишком маленькое \(\delta\) не исправит light bleeding
\item Слишком большое \(\delta\) сделает тени слишком жёсткими
\item Нужно подобрать самое маленькое подходящее значение :)
\end{itemize}
\end{frame}

\begin{frame}[fragile]
\frametitle{Задание 5}
\slideimage{5.png}
\end{frame}

\begin{frame}[fragile]
\frametitle{Задание 6*}
Размываем тени
\begin{itemize}
\item Перед вычислением \(\mu\) и \(\sigma\), усредняем значения из shadow map с гауссовыми весами
\item \textbf{NB:} Усреднять нужно именно \textbf{значения из самой текстуры}, а не вычисленные через неравенство Чебышева значения освещённости
\end{itemize}
\end{frame}

\begin{frame}[fragile]
\frametitle{Задание 6*}
\slideimage{6.png}
\end{frame}

\end{document}
